\documentclass[11pt]{article}
\usepackage[T1]{fontenc}
\usepackage{lmodern}
\usepackage[margin=1in]{geometry}
\usepackage{amsmath,amssymb,amsthm}
\usepackage{microtype}
\usepackage[colorlinks=true,linkcolor=blue,citecolor=blue,urlcolor=blue]{hyperref}
\hypersetup{pdftitle={Rational Hodge classes on mixed products of K3 moduli spaces},pdfauthor={Alec Kriebel},pdfsubject={Rational Hodge conjecture; K3 surfaces; moduli spaces; Chow motives}}
\setlength{\emergencystretch}{2em}
\newtheorem{theorem}{Theorem}
\newtheorem{proposition}[theorem]{Proposition}
\newtheorem{remark}[theorem]{Remark}
\newcommand{\Q}{\mathbb Q}
\newcommand{\CH}{\operatorname{CH}}
\newcommand{\Hdg}{\operatorname{Hdg}}
\newcommand{\cl}{\operatorname{cl}}
\newcommand{\Br}{\operatorname{Br}}
\title{Rational Hodge classes on mixed products\newline of K3 moduli spaces}
\author{Alec Kriebel\thanks{ORCID: \href{https://orcid.org/0009-0001-9320-500X}{0009-0001-9320-500X}.}}
\date{October 6, 2026}
\begin{document}
\maketitle
\begin{abstract}
Combining OpenAI's rational Hodge theorem for arbitrary mixed products of
projective K3 surfaces with B\"ulles' Chow-motive decomposition, we record
rational Hodge in every codimension for finite mixed products of the smooth
projective K3 moduli spaces covered by that decomposition. The factors may
have different K3 bases and include Hilbert schemes, twisted stable-sheaf
moduli, and eligible moduli of Bridgeland-stable objects. We give the explicit
pull-push argument, retaining smoothness, projectivity, and stability
hypotheses. The contribution is an assembled corollary of the cited results;
it supplies neither a new proof of the base K3 theorem nor a new motivic
reduction.
\end{abstract}

\section{Statement and provenance}
For a smooth projective complex variety $X$, write
\[
 \Hdg^p(X)=H^{2p}(X,\Q)\cap H^{p,p}(X,\mathbb C).
\]
Rational Hodge in codimension $p$ means surjectivity of
$\cl:\CH^p(X)_\Q\to\Hdg^p(X)$. Here Chow groups and all correspondences
have rational coefficients.

\begin{theorem}\label{thm:moduli}
Let $S_1,\ldots,S_r$ be complex projective K3 surfaces and let
$\alpha_i\in\Br(S_i)$. For each $i$, let $M_i$ be a smooth projective
moduli space of Gieseker-stable $\alpha_i$-twisted sheaves, or a smooth
projective moduli space of $\sigma_i$-stable objects in
$D^b(S_i,\alpha_i)$ for a generic Bridgeland stability condition
$\sigma_i$, as in B\"ulles' Theorem~0.1 \cite{Bulles}.
Then the rational Hodge conjecture holds in every codimension on
\[
 X=\prod_{i=1}^r M_i.
\]
The K3 bases may be different or repeated. In particular, the conclusion
holds for every finite mixed product of Hilbert schemes $S_i^{[n_i]}$,
$n_i\geq0$.
\end{theorem}

The sheaf data and polarization defining each Gieseker moduli problem
are fixed. Genericity for Bridgeland stability is relative to the fixed
numerical data of that moduli problem. Smoothness and projectivity of the
actual stable moduli space are hypotheses, rather than conclusions about
arbitrary numerical data or stability conditions. The theorem does not
add a primitivity or fine-moduli assumption to B\"ulles' statement.
Its exact stable-object scope is that of the cited theorem.

\begin{samepage}
The principal input is OpenAI's theorem \cite[Theorem~1.1]{OpenAIMixed}
that rational Hodge holds on every finite product of projective complex
K3 surfaces, with arbitrary distinct or repeated factors. We use the
public OpenAI Math Release snapshot at the following commit, accessed
October~6, 2026:
\begin{center}
\texttt{adc7f1241b42e322a6451854ab7e4b4c146bf78a}.
\end{center}
Manuscript cover dates in the references are not assertions about
their earliest public disclosure.
\end{samepage}

Arapura had already proved conditional Hodge transfers to Hilbert schemes
and K3 sheaf moduli \cite[Theorems~5.4 and~5.7(3)]{ArapuraHodge}, and a
broader transfer for motivated varieties \cite[Lemma~4.2 and
Theorems~7.4, 7.8]{ArapuraMotivation}. B\"ulles supplies the precise
unconditional Chow-motive splitting for twisted sheaves and the eligible
Bridgeland moduli \cite[Theorem~0.1]{Bulles}. The earlier OpenAI
quadratic-locus companion already records the moduli and self-power
application on its restricted locus \cite[Corollary~9.5]{OpenAIQuadratic}.
Its general exact Kuga--Satake criterion, together with the later universal
Kuga--Satake theorem \cite{OpenAIKS}, also yields the unrestricted-base
self-power application. We record the mixed-moduli assembly: the extra
mixed-base quantifier is supplied by \cite{OpenAIMixed}, and the reduction
itself is inherited. We make no claim of a first proof or substantial
independent novelty.

\section{The transfer with degrees made explicit}

For a smooth projective $V$ of dimension $d_V$ and a smooth projective
$W$, a correspondence $A\in\CH^a(V\times W)_\Q$ acts by
\[
 A_*z=(\operatorname{pr}_W)_*
       \bigl(\operatorname{pr}_V^*z\cdot A\bigr).
\]
On cohomology this sends a rational $(p,p)$ class to a rational
$(p+a-d_V,p+a-d_V)$ class. Cycle classes commute with this action.
If $B\in\CH^b(W\times V)_\Q$, then $B\circ A$ has codimension
$a+b-\dim W$ in $V\times V$.

\begin{proposition}\label{prop:transfer}
Suppose a smooth projective $X$ is a direct summand of a finite sum of
Tate-twisted rational Chow motives of smooth projective varieties $K_j$.
If rational Hodge holds in every codimension on each $K_j$, it holds in
every codimension on $X$.
\end{proposition}

\begin{proof}
The split maps and projectors give homogeneous rational algebraic
correspondences
\[
 A_j\in\CH^{a_j}(X\times K_j)_\Q,\qquad
 B_j\in\CH^{b_j}(K_j\times X)_\Q,
\]
with
\begin{equation}\label{eq:splitting}
 a_j+b_j=\dim X+\dim K_j,\qquad
 \sum_j B_j\circ A_j=\Delta_X.
\end{equation}
Using the actual cycle codimensions avoids any convention-dependent
sign for a Tate twist.
Let $z\in\Hdg^p(X)$ and put $q_j=p+a_j-\dim X$. The class $A_{j*}z$
lies in $\Hdg^{q_j}(K_j)$. If $q_j$ is in the cohomological range,
choose $Z_j\in\CH^{q_j}(K_j)_\Q$ with
$\cl(Z_j)=A_{j*}z$ by the assumed Hodge conjecture. Outside that range
the target is zero, and take $Z_j=0$. Then
\[
 Z=\sum_j B_{j*}Z_j\in\CH^p(X)_\Q,
 \qquad
 \cl(Z)=\sum_j B_{j*}A_{j*}z=z.
\]
The codimension is $q_j+b_j-\dim K_j=p$ by
\eqref{eq:splitting}, and the final equality uses the diagonal identity.
\end{proof}

\begin{proof}[Proof of Theorem~\ref{thm:moduli}]
First assume that all factors are nonempty and positive dimensional.
B\"ulles gives split maps
\[
 h(M_i)\longrightarrow
   \bigoplus_{j\in J_i}h(S_i^{k_{ij}})(t_{ij})
   \longrightarrow h(M_i)
\]
whose composite is the identity, with finite $J_i$,
$1\leq k_{ij}\leq\dim M_i$, and integers $t_{ij}$. Tensor these maps
in rational Chow motives and distribute the finite sums. The motive of
$X$ is a summand of a finite sum of Tate twists of motives of
\[
 K_J=\prod_{i=1}^r S_i^{k_{ij_i}},
 \qquad J=(j_1,\ldots,j_r)\in\prod_i J_i.
\]
At the level of correspondences these are exterior products of the
individual split maps; their composite sum is the exterior product of
the diagonals, hence $\Delta_X$. Their codimensions add, so they satisfy
\eqref{eq:splitting}. The rational Hodge theorem for mixed K3 products
\cite{OpenAIMixed} applies to every $K_J$. Proposition~\ref{prop:transfer}
therefore proves the claim.

An empty factor gives an empty product. A smooth projective
zero-dimensional factor over $\mathbb C$ is a finite disjoint union of
points; remove it on each component and apply the preceding argument.
The empty product is a point. Codimensions outside $[0,\dim X]$ have
zero target.
For $n\geq1$, $S^{[n]}$ is the smooth projective moduli of rank-one ideal
sheaves with Mukai vector $(1,0,1-n)$ and is in the untwisted sheaf case.
For $n=0$ it is a point. This proves the stated Hilbert-scheme consequence.
\end{proof}

\begin{remark}[Non-fine and twisted moduli]
The split maps used here are B\"ulles' algebraic Chow correspondences.
We do not construct them by presuming that an honest universal family
exists. The cited construction uses quasi-universal/twisted data and the
normalized diagonal identity underlying Markman's result
\cite[Section~3]{Markman}. Raw duplication of a universal family does
not preserve its Ext diagonal formula: multiplicities must be normalized.
The present argument needs only the actual rational motive splitting,
which already covers the Brauer classes in Theorem~\ref{thm:moduli}.
\end{remark}

\section{Scope and verification}
The transfer explicitly uses Hodge for mixed products $K_J$; separate
self-power theorems alone do not supply this premise. Conversely,
$S^{[1]}=S$ shows that the unrestricted mixed-moduli conclusion includes
the mixed-K3 base theorem. Thus, given B\"ulles' established reduction,
these two universal assertions are equivalent. This explains exactly
where the principal mathematical input lies.

No claim is made about arbitrary deformations of $K3^{[n]}$ type,
singular moduli or semistable compactifications, moduli on arbitrary
abelian surfaces, the integral or generalized Hodge conjecture, or
finite-dimensionality of Chow motives.

The verification materials distinguish the explicit correspondence
argument from the unformalized geometric input in \cite{OpenAIMixed}.
The upstream Kuga--Satake Lean module checks rational Clifford/tensor
identities; it does not formalize algebraic cycles, the universal
Kuga--Satake theorem, or the theorem of this note. Reproducible small
linear-algebra and sign checks are supporting evidence for parts of the
upstream audit, not certificates for its full geometry.

\paragraph{AI use and review status.}
AI tools were used extensively in research, source and priority checks,
drafting, and verification. The cited OpenAI manuscripts are attributed
to OpenAI according to their supplied citation information. This note has
not undergone conventional human peer review or refereeing. Automated
adversarial reviews are not human peer review.

\begin{thebibliography}{9}
\raggedright
\bibitem{ArapuraHodge}
D. Arapura, \emph{Hodge cycles on some moduli spaces},
\href{https://arxiv.org/abs/math/0102070v5}{arXiv:math/0102070v5}, 2002.
\bibitem{ArapuraMotivation}
D. Arapura, \emph{Motivation for Hodge cycles}, Adv. Math. \textbf{207}
(2006), 762--781, \href{https://doi.org/10.1016/j.aim.2006.01.005}{doi:10.1016/j.aim.2006.01.005}.
\bibitem{Bulles}
T.-H. B\"ulles, \emph{Motives of moduli spaces on K3 surfaces and of special
cubic fourfolds}, Manuscripta Math. \textbf{161} (2020), 109--124,
\href{https://arxiv.org/abs/1806.08284v1}{arXiv:1806.08284v1},
\href{https://doi.org/10.1007/s00229-018-1086-0}{doi:10.1007/s00229-018-1086-0}.
\bibitem{Markman}
E. Markman, \emph{Generators of the cohomology ring of moduli spaces of
sheaves on symplectic surfaces}, J. reine angew. Math. \textbf{544}
(2002), 61--82, \href{https://arxiv.org/abs/math/0009109v4}{arXiv:math/0009109v4}.
\bibitem{OpenAIQuadratic}
OpenAI, \emph{Algebraic Kuga--Satake correspondences and Hodge conjectures
on a K3 quadratic locus}, OpenAI Math Release preprint, cover date
September 30, 2026, \href{https://github.com/openai/math/blob/adc7f1241b42e322a6451854ab7e4b4c146bf78a/preprints/Algebraic-Kuga-Satake-correspondences-and-Hodge-conjectures-on-a-K3-quadratic-locus-September-30-2026/paper.pdf}{pinned public manuscript}.
\bibitem{OpenAIKS}
OpenAI, \emph{Algebraicity of Kuga--Satake Correspondences for K3 Surfaces},
OpenAI Math Release preprint, cover date October 3, 2026,
\href{https://github.com/openai/math/blob/adc7f1241b42e322a6451854ab7e4b4c146bf78a/preprints/Algebraicity-of-Kuga-Satake-Correspondences-for-K3-Surfaces-October-3-2026/manuscript.pdf}{pinned public manuscript}.
\bibitem{OpenAIMixed}
OpenAI, \emph{The rational Hodge conjecture for products of K3 surfaces},
OpenAI Math Release preprint, cover date October 4, 2026,
\href{https://github.com/openai/math/blob/adc7f1241b42e322a6451854ab7e4b4c146bf78a/preprints/The-rational-Hodge-conjecture-for-products-of-K3-surfaces-October-4-2026/hodge-conjecture-products-k3.pdf}{pinned public manuscript}.
\end{thebibliography}
\end{document}
